\chapter{\label{ch:laniga}Model of LaNiGa\textsubscript{2}}
\graphicspath{{../gfx/chapter03/}{../plots/chapter03/}}


In this section we propose a tight-binding model of LaNiGa$_2$. We model only the interactions between Nickel atoms, as a minimal model of superconductivity. The Nickel atoms form a diatomic lattice, with two orbitals per site. We model in the plane such that the atoms are spatially degenerate. Therefore, the model is described by an onsite chemical potential tensor
\begin{equation}
\bm{\mu}\equiv\bm{\mu}_\text{sites}\otimes\bm{\mu}_\text{spin}\otimes\bm{\mu}_\text{orbitals},
\end{equation}
where
\begin{equation}
\bm{\mu}_\text{sites}\equiv
\begin{pmatrix}
-\mu_a && 0 \\ 0 && -\mu_b
\end{pmatrix},
\end{equation}
\begin{equation}
\bm{\mu}_\text{spin}\equiv\mathbb{1}_2,
\end{equation}
\begin{equation}
\bm{\mu}_\text{orbitals}\equiv
\begin{pmatrix}
1+s && \delta \\ \delta && 1-s
\end{pmatrix},
\end{equation}
and where $\mu_a$ and $\mu_b$ refer to the onsite chemical potentials for the atoms \textit{above} and \textit{below} the plane. The orbitals include a rigid shift $s$ and hybridisation factor $\delta$.

The intercell nearest-neighbour hopping is described by the tensor
\begin{equation}
\bm{T}\equiv\bm{T}_\text{sites}\otimes\bm{T}_\text{spin}\otimes\bm{T}_\text{orbitals},
\end{equation}
where
\begin{equation}
\bm{T}_\text{sites}\equiv
\begin{pmatrix}
-t && -t' \\ -t' && -t
\end{pmatrix},
\end{equation}
\begin{equation}
\bm{T}_\text{spin}\equiv\mathbb{1}_2,
\end{equation}
\begin{equation}
\bm{T}_\text{orbitals}\equiv
\begin{pmatrix}
1+\alpha && 1+\beta \\ 1+\beta && 1-\alpha
\end{pmatrix},
\end{equation}
and where $t>0$ refers to intrasite and $t'<t$ refers to intersite tunnelling processes. $\alpha>0$ describes intraorbital hopping ($d_{xy}-d_{xy}$, $d_{z^2}-d_{z^2}$) and $\beta>0$ describes interorbital orbital hopping ($d_{xy}-d_{z^2}$, $d_{z^2}-d_{xy}$).

